\pagebreak

\chapter{regular expressions}\label{chap:regexp}

    The following summary of regular expressions is taken directly from the Java
    documentation -

    \begin{verbatim}
    Characters

        x  The character x
        \\  The backslash character
        \0n  The character with octal value 0n  (0 <= n <= 7)
        \0nn  The character with octal value 0nn  (0 <= n <= 7)
        \0mnn  The character with octal value 0mnn  (0 <= m <= 3, 0 <= n <= 7)
        \xhh  The character with hexadecimal value 0xhh
        \uhhhh  The character with hexadecimal value 0xhhhh
        \t  The tab character ('\u0009')
        \n  The newline (line feed) character ('\u000A')
        \r  The carriage-return character ('\u000D')
        \f  The form-feed character ('\u000C')
        \a  The alert (bell) character ('\u0007')
        \e  The escape character ('\u001B')
        \cx  The control character corresponding to x
    \end{verbatim}

    \begin{verbatim}
    Character classes

        [abc]  a, b, or c (simple class)
        [^abc]  Any character except a, b, or c (negation)
        [a-zA-Z]  a through z  or A through Z, inclusive (range)
        [a-d[m-p]]  a through d, or m through p:
        [a-dm-p] (union)  [a-z&&[def]]  d, e, or f (intersection)
        [a-z&&[^bc]]  a through z, except for b and c:
        [ad-z] (subtraction)
        [a-z&&[^m-p]]  a through z, and not m through p:
        [a-lq-z](subtraction)
    \end{verbatim}

    \begin{verbatim}
    Predefined character classes

        .  Any character (may or may not match line terminators)
        \d  A digit: [0-9]
        \D  A non-digit: [^0-9]
        \s  A whitespace character: [ \t\n\x0B\f\r]
        \S  A non-whitespace character: [^\s]
        \w  A word character: [a-zA-Z_0-9]
        \W  A non-word character: [^\w]     
    \end{verbatim}

    \begin{verbatim}
    Boundary matchers
    ^	The beginning of a line
    $	The end of a line
    \b	A word boundary
    \B	A non-word boundary
    \A	The beginning of the input
    \G	The end of the previous match
    \Z	The end of the input but for the final terminator, if any
    \z	The end of the input
    \end{verbatim}

    \begin{verbatim}
    Greedy quantifiers
    X?	X, once or not at all
    X*	X, zero or more times
    X+	X, one or more times
    X{n}	X, exactly n times
    X{n,}	X, at least n times
    X{n,m}	X, at least n but not more than m times
    \end{verbatim}

    \begin{verbatim}
    Reluctant quantifiers
    X??	X, once or not at all
    X*?	X, zero or more times
    X+?	X, one or more times
    X{n}?	X, exactly n times
    X{n,}?	X, at least n times
    X{n,m}?	X, at least n but not more than m times
    \end{verbatim}

    \begin{verbatim}
    Possessive quantifiers
    X?+	X, once or not at all
    X*+	X, zero or more times
    X++	X, one or more times
    X{n}+	X, exactly n times
    X{n,}+	X, at least n times
    X{n,m}+	X, at least n but not more than m times
    \end{verbatim}

    \begin{verbatim}
    Logical operators
    XY	X followed by Y
    X|Y	Either X or Y
    (X)	X, as a capturing group
    \end{verbatim}

    \begin{verbatim}
    Back references
    \n	Whatever the nth capturing group matched
    \end{verbatim}

    \begin{verbatim}
    Quotation
    \	Nothing, but quotes the following character
    \Q	Nothing, but quotes all characters until \E
    \E	Nothing, but ends quoting started by \Q
    \end{verbatim}

    \begin{verbatim}
    Special constructs (non-capturing)
    (?:X)	X, as a non-capturing group
    (?idmsux-idmsux) 	Nothing, but turns match flags i d m s u x on - off
    (?idmsux-idmsux:X)  	X, as a non-capturing group with the given flags i d m s u x on - off
    (?=X)	X, via zero-width positive lookahead
    (?!X)	X, via zero-width negative lookahead
    (?<=X)	X, via zero-width positive lookbehind
    (?<!X)	X, via zero-width negative lookbehind
    (?>X)	X, as an independent, non-capturing group
    Backslashes, escapes, and quoting
    \end{verbatim}

    \begin{verbatim}
    The backslash character ('\') serves to introduce escaped constructs, as
    defined in the table above, as well as to quote characters that otherwise
    would be interpreted as unescaped constructs. Thus the expression \\ matches
    a single backslash and \{ matches a left brace.

    It is an error to use a backslash prior to any alphabetic character that
    does not denote an escaped construct; these are reserved for future
    extensions to the regular-expression language. A backslash may be used prior
    to a non-alphabetic character regardless of whether that character is part
    of an unescaped construct.

    Backslashes within string literals in Java source code are interpreted as
    required by the Java Language Specification as either Unicode escapes or
    other character escapes. It is therefore necessary to double backslashes in
    string literals that represent regular expressions to protect them from
    interpretation by the Java bytecode compiler. The string literal "\b", for
    example, matches a single backspace character when interpreted as a regular
    expression, while "\\b" matches a word boundary. The string literal
    "\(hello\)" is illegal and leads to a compile-time error; in order to match
    the string (hello) the string literal "\\(hello\\)" must be used.
    \end{verbatim}











    \begin{verbatim}
    Groups and capturing

    Capturing groups are numbered by counting their opening parentheses from
    left to right. In the expression ((A)(B(C))), for example, there are four
    such groups:

    1    	((A)(B(C)))
    2    	(A)
    3    	(B(C))
    4    	(C)
    Group zero always stands for the entire expression.

    Capturing groups are so named because, during a match, each subsequence of
    the input sequence that matches such a group is saved. The captured
    subsequence may be used later in the expression, via a back reference, and
    may also be retrieved from the matcher once the match operation is complete.

    The captured input associated with a group is always the subsequence that
    the group most recently matched. If a group is evaluated a second time
    because of quantification then its previously-captured value, if any, will
    be retained if the second evaluation fails. Matching the string "aba"
    against the expression (a(b)?)+, for example, leaves group two set to "b".
    All captured input is discarded at the beginning of each match.

    Groups beginning with (? are pure, non-capturing groups that do not capture
    text and do not count towards the group total.
    \end{verbatim}
